\documentclass[utf8,usehyperref,12pt]{G7-32}
\input glyphtounicode.tex                   % Нужно для
\pdfgentounicode=1                          % поиска по pdf в кириллице
\usepackage[T2A]{fontenc}
\usepackage{pscyr}
\usepackage[utf8]{inputenc}                 % ваша любимая кодировка здесь
\usepackage[english,russian]{babel}         % это необходимо для включения переносов
\usepackage{float}
\usepackage[pdftex]{graphicx}
\usepackage{array}
\usepackage{mversion}                       % Версионирование
\TableInChaper                              % таблицы будут нумероваться в пределах раздела
\PicInChaper                                % рисунки будут нумероваться в пределах раздела
\setlength\GostItemGap{2mm}                 % для красоты можно менять от 0мм
\sloppy                                     % переносы

\newcounter{MinoreCounter}                  % Новый счетчик минорных версий
\setcounter{MinoreCounter}{1}               % Текущая минорная версия

\renewcommand{\theequation}{(\arabic{chapter}.\arabic{equation})}   % Формат нумерации формул "Глава.номер"
\renewcommand{\version}{\versionnumber.\theMinoreCounter.\thebuildcounter} % Формат версии

\makeatletter
\@addtoreset{equation}{chapter}             % Счетчик формул
\makeatother

\setcounter{page}{3}                        % Начало нумерации страниц
\setVersion{0}                              % Текущая мажорная версия
\increaseBuild                              % Увеличивать номер сборки при каждой компиляции



  
    

%<------------- НАЧАЛО ДОКУМЕНТА ---------------------------

\begin{document}
\tableofcontents

\begin{center}
Текущая версия документа: \version~от \today 
\end{center}

% Это самый главный файл документа, тут лучше ничего не писать, а только собирать
\chapter{Название главы}

\begin{itemize}
	\item пункт 1
	\item пункт 2
	\item пункт 3
	\begin{itemize}
		\item подпункт 1
		\item подпункт 1
		\item подпункт 1
	\end{itemize}
	\item пункт 4
\end{itemize}

\section{Секция}
\subsection{Подсекция}
\subsubsection{ПодПодсекция}

\begin{figure}[H]
\begin{center}
\includegraphics[scale=0.6]{./image/1.jpg}
\end{center}
\caption{Картина}
\end{figure}

\textit{курсив}, \textbf{жирный}\footnote{сноска} , \verb\моноширный\, 

Lorem ipsum dolor sit amet, consectetur adipiscing elit. Cras cursus neque vel libero iaculis vehicula. 
продолжаем текст на этой же строке

Что бы сделать новый абзац надо дважды поставить символ переноса строки(Enter) Lorem ipsum dolor sit amet, consectetur adipiscing elit. 

Вот так, опять новый абзац Lorem ipsum dolor sit amet, cursus consectetur venenatis. Fusce blandit orci a neque ultrices, id vestibulum risus scelerisque. Sed at magna odio.

Если надо вставить отдельно какой то код то делаем так:
\begin{verbatim}
class Mid{
/*
 * Какой то комментарий
 */
public:
	// Константы при работе с группами портов
	static const unsigned int GROUP_ALL = 0x0;
	static const unsigned int GROUP_A = 0x1;
	static const unsigned int GROUP_B = 0x2;
	static const unsigned int GROUP_C = 0x3;
	static const unsigned int GROUP_D1 = 0x4;
\end{verbatim}

Теперь формулы, если надо вставить формулу с нумерацией, то тогда так:
\begin{equation}\label{Example} %это название ссылки на эту формулу
k^{10^{-6}} = \sqrt{\frac{Q_{i}}{10^{12}}}
\end{equation}
или мы хотим в тексте вставить формулу без номера $ \forall x \in \Re ~\exists y, ~y_{i} \neq \int \sum^{0}_{i} $

Да, знак \verb\~\ означает неразрывный пробел, а если надо ввести какой то служебный символ латеха, чтобы он не обрабатывался, ставится обратный слеш, например \%

Если надо сослаться на формулу \ref{Example}

И на последок, что бы автоматически сгенерировалось оглавление и ссылки, нужно дважды скомпилировать документ, через F1 или кнопку PDFLaTeX в Texmaker 
%\chapter{Теоретические сведения}

\section{Основы программирования на языке \textit{Python}}

\section{Основы работы с системой \textit{GitHub}}

\section{Основы работы с текстовым редактором \LaTeX{}}

\section{Справочные материалы по радиотехнике}

\subsection{Источник информации}

\textbf{Источник информации} --- устройство, передающее информацию по средствам цифровой системы связи. Источник информации может быть аналоговым и дискретным. Выход аналогового источника может иметь любое значение из непрерывного диапазона амплитуд, тогда как выход источника дискретной информации --- значения из конечного множества амплитуд. Источники аналоговой информации преобразуются в источники цифровой информации по средствам квантования и дискретизации.

В лабораторных работах будет использоваться только дискретный, бинарный источник информации с равновероятным возникновением одного из двух событий (выпадение 0 или 1).

\subsection{Каналы связи, шумы, помехи}

\textbf{Канал связи} --- среда распространения, или электромагнитный тракт, соединяющий передатчик и приемник. Каналы связи могут быть воздушные, кабельные, световодные, тропосферные, спутниковые, непрерывные (на входе и выходе канала - непрерывные сигналы), дискретные или цифровые (на входе и выходе канала - дискретные сигналы), непрерывно-дискретные (на входе канала - непрерывные сигналы, а на выходе - дискретные сигналы), дискретно-непрерывные (на входе канала - дискретные сигналы, а на выходе - непрерывные сигналы) и др. В любом канале связи к сигналу добавляется шум или помеха.

Наиболее распространенная модель шума, используемая в теории связи --- аддитивный белый гауссовский шум (АБГШ). Из названия следует, что АБГШ складывается с сигналом в канале связи. Такой шум с точки зрения статистики имеет гауссовское распределение.

\subsection{Помехоустойчивое кодирование}
\textbf{Помехоустойчивое кодирование} вносит избыточность в исходную информационную последовательность. На основе дополнительной (избыточной) информации появляется возможность обнаруживать и исправлять ошибки с определенной кратностью. Кратности исправляемой и определяемой ошибок определяются свойствами конкретного помехоустойчивого кода, в том числе количеством избыточных символов, добавляемых к исходной последовательности.
\subsection{Модуляция}

\subsection{Статистические характеристики системы} % В разработке
%\chapter{Пример выполнения задания}

\section{Вариант задания-примера}

\section{Текст программы на языке \textit{Python}}

\section{Сопроводительная документация}


\chapter{Варианты заданий}
Во всех вариантах использовать дискретный бинарный источник информации (равновероятное возникновение нулей и единиц). Длина последовательности 512 символов. Модуль сжатия исходной информации используется на усмотрение студентов (использование не обязательно). В канале связи присутствует только аддитивный белый гауссовский шум. Мощность сигнала и шума выбирать таким образом, чтобы отношение сигнал/шум  Остальные требования указаны в таблице.


\begin{center}
\begin{tabular}{|c|p{4cm}|c|c|}
\hline
Номер варианта & Методы кодирования & Кратность исправляемой ошибки & Виды модуляции \\
\hline
1 & код Рида-Саломона и БЧХ-код & 5 & QAM-128 и PM-16\\
\hline
2 & код Хемминга и БЧХ-код & 7 & QAM-256 и QPSK \\
\hline
\end{tabular}
\end{center}

 

\end{document}